\section{An exact relation for cyclic deformations}
\label{sec:uniform-torsion}

An explicit exact sequence annihilates a cyclic deformation class by
the truncation exponent.  The next section constructs an invariant
that proves this bound is exact for every exponent.

\begin{proposition}\label{prop:truncation-torsion}
Let $p$ be a prime, $O=\Z_{(p)}$ or $\Z_p$, $n\geq2$, and
$A=O[e]/(e^n)$.
Let $C$ be the companion of $\Phi_p(z)$ in the basis
$1,z,\ldots,z^{p-2}$, let $a=e_{p-1}^{\mathsf T}$ and $b=e_1$, and put
\[
 T_j=\begin{pmatrix}1&e^j a\\e^{n-j}b&C\end{pmatrix},
 \qquad L_j=(A^p,T_j),\qquad 0\leq j\leq n.
\]
If $x_j=[L_j]-[A[C_p]]$, then $x_j=jx_1$ and $nx_1=0$.
In particular, the flat dual-number obstruction of
Theorem~\ref{thm:dual-numbers-all-primes} has exact order two for every
prime $p$, both before and after induction to $S_p$.
\end{proposition}

\begin{proof}
The zero cross products and $\Phi_p(C)=0$ give $T_j^p=I$.
The endpoint lattices are regular.  Indeed, multiplication by a
generator $g$ on $A[C_p]$ has matrix $T_0$ in the basis
\[
 1+g+\cdots+g^{p-1},\quad 1,g,\ldots,g^{p-2},
\]
and matrix $T_n$ in the basis
\[
 1,\quad g-1,(g-1)g,\ldots,(g-1)g^{p-2}.
\]
Both lists are bases by triangular comparison with the powers of $g$.

For $0\leq j<n$, define an equivariant injection
\[
 L_j\longrightarrow L_{j+1}\oplus L_0,\qquad
 (v_0,v_c)\longmapsto
 \bigl((ev_0,v_c),(v_0,e^jv_c)\bigr).
\]
Its two blocks are $\operatorname{diag}(e,I_{p-1})$ and
$\operatorname{diag}(1,e^jI_{p-1})$; direct multiplication verifies
equivariance.  It is split over $A$.  Quotient coordinates are
\[
 q_0=x_0-ey_0,\qquad q_c=y_c-e^jx_c,
\]
where $x$ and $y$ denote the two target summands.  In these coordinates
the induced action is
\[
 \begin{pmatrix}1&-e a\\-e^{n-1}b&C\end{pmatrix},
\]
which is conjugate to $T_1$ by $\operatorname{diag}(1,-I_{p-1})$.
Thus there is an admissible exact sequence
\[
 0\longrightarrow L_j\longrightarrow L_{j+1}\oplus L_0
   \longrightarrow L_1\longrightarrow0.
\]
It gives $x_{j+1}=x_j+x_1$.  Since $x_0=x_n=0$, the two assertions
follow.  For $n=2$, the previously constructed detector is nonzero
on $x_1$ and on its induction to $S_p$, whereas the displayed relation
kills both classes by two.  They therefore have exact order two.
\end{proof}

Theorem~\ref{thm:all-truncations} proves that this annihilation bound
is exact, both for the cyclic classes and after induction.
